\chapter{Frobenius 定理：从分布到积分叶}

向量场给出“允许运动的方向”，但给定若干方向并不保证它们真能拼成低维子流形。例如在三维空间中，每点选一个二维平面场，这些平面可能像螺旋一样彼此扭转，因而不存在处处与它们相切的曲面。Frobenius 定理精确回答可积性问题：一个常秩分布局部来自子流形，当且仅当它在 Lie 括号下闭合；用余切语言说，当且仅当它的湮灭一形式生成的微分理想在 \(\dd\) 下闭合。

\section{分布、积分流形与 involutivity}

秩为 \(k\) 的光滑分布是切丛的秩 \(k\) 光滑子丛
\[
\mathcal D\subset TM.
\]
局部可以选光滑向量场 \(X_1,\ldots,X_k\) 使
\[
\mathcal D_p
=\operatorname{span}\{X_1(p),\ldots,X_k(p)\}.
\]
若存在浸入的 \(k\) 维子流形 \(L\subset M\)，满足每点
\[
T_pL=\mathcal D_p,
\qquad p\in L,
\]
则称 \(L\) 为 \(\mathcal D\) 的积分流形。若每一点附近都能由这样的积分流形形成局部叶层，则称 \(\mathcal D\) 可积。

分布称为 involutive，若任意局部截面 \(X,Y\in\Gamma(\mathcal D)\) 都满足
\begin{equation}\label{eq:math-dg-frobenius-involutive}
[X,Y]\in\Gamma(\mathcal D).
\end{equation}
这个条件完全由光滑结构决定，不需要联络或度量。

\begin{theorem}[Frobenius 定理]\label{thm:math-dg-frobenius}
设 \(\mathcal D\subset TM\) 是常秩 \(k\) 的光滑分布。以下条件在局部等价：
\begin{enumerate}[label=(\roman*)]
\item \(\mathcal D\) 可积；
\item \(\mathcal D\) involutive，即满足\eqnrefs{eq:math-dg-frobenius-involutive}；
\item 每点附近存在坐标 \((y^1,\ldots,y^n)\)，使
\[
\mathcal D
=\operatorname{span}\{\partial_{y^1},\ldots,\partial_{y^k}\};
\]
\item 若 \(\alpha^1,\ldots,\alpha^{n-k}\) 是湮灭丛
\[
\mathcal D^\circ
:=\{\alpha\in T^*M:\alpha|_{\mathcal D}=0\}
\]
的一组局部标架，则存在一形式 \(\eta^a{}_b\) 使
\begin{equation}\label{eq:math-dg-frobenius-pfaff}
\dd\alpha^a=\eta^a{}_b\wedge\alpha^b.
\end{equation}
\end{enumerate}
\end{theorem}

从 (iii) 到 (i) 是直接的：固定 \(y^{k+1},\ldots,y^n\) 即得到积分叶。从 (i) 到 (ii) 也立即成立，因为切于同一子流形的向量场的 Lie 括号仍切于该子流形。真正困难的是由括号闭合构造适配坐标。

\begin{derivation}[由 involutivity 构造局部积分叶]
对秩 \(k\) 归纳。\(k=1\) 时，取非零局部截面 \(X_1\)。flow-box 定理给出坐标使
\[
X_1=\partial_{y^1},
\]
故积分曲线就是一维叶。

设结论对秩 \(k-1\) 成立。取 \(p\) 附近非零截面 \(X_1\in\Gamma(\mathcal D)\)，再次用 flow-box 坐标令
\[
X_1=\partial_{y^1}.
\]
令横截面
\[
\Sigma:=\{y^1=0\}.
\]
通过对其余局部生成元减去适当的 \(X_1\) 分量，可以选取
\[
Y_2,\ldots,Y_k\in\Gamma(\mathcal D)
\]
使它们在 \(\Sigma\) 上切于 \(\Sigma\)。于是
\[
\mathcal E:=\mathcal D\cap T\Sigma
\]
是 \(\Sigma\) 上秩 \(k-1\) 分布。若 \(Y,Z\) 是 \(\mathcal E\) 的截面，则它们既切于 \(\Sigma\) 又属于 \(\mathcal D\)。切于固定子流形的向量场在 Lie 括号下封闭，而 \(\mathcal D\) 假设 involutive，因此
\[
[Y,Z]\in\Gamma(\mathcal D\cap T\Sigma)=\Gamma(\mathcal E).
\]
由归纳假设，缩小邻域后可在 \(\Sigma\) 上取适配坐标
\[
(y^2,\ldots,y^k,y^{k+1},\ldots,y^n),
\]
使 \(\mathcal E\) 由前 \(k-1\) 个坐标方向张成；固定
\(y^{k+1},\ldots,y^n\) 得到 \(\mathcal E\) 的局部积分叶族 \(L_c\)。特别地，其中一片 \(L_0\) 通过 \(p\)。

现在令 \(\Phi_t\) 为 \(X_1\) 的局部流。关键是证明它保持分布 \(\mathcal D\)。取局部标架 \(X_1,\ldots,X_k\) 张成 \(\mathcal D\)。involutive 条件给出
\[
[X_1,X_a]=B^b{}_aX_b
\]
其中 \(B^b{}_a\) 为光滑函数。沿 \(X_1\) 的积分曲线，推前标架满足线性常微分方程；等价地，\((\Phi_{-t})_*X_a\) 的导数仍是这些向量的线性组合。因此初值位于 \(\mathcal D\) 的向量在整个局部流中始终留在 \(\mathcal D\)，即
\[
(\Phi_t)_*\mathcal D=\mathcal D.
\]

对 \(\Sigma\) 中每一片局部叶 \(L_c\) 定义
\[
L_c^{\mathrm{sat}}
:=\{\Phi_t(q):q\in L_c,\ |t|<\varepsilon\}.
\]
其切空间由 \(X_1\) 与 \((\Phi_t)_*T_qL_c\) 张成；前者属于 \(\mathcal D\)，后者由于分布被流保持也属于 \(\mathcal D\)。维数同为 \(k\)，故
\[
T_{\Phi_t(q)}L_c^{\mathrm{sat}}
=\mathcal D_{\Phi_t(q)}.
\]
因此这些饱和叶给出 \(p\) 附近的局部叶层。把流参数作为 \(y^1\)，保留 \(\Sigma\) 上的
\(y^2,\ldots,y^n\)，便得到局部坐标，其中
\[
\mathcal D
=\operatorname{span}\{\partial_{y^1},\ldots,\partial_{y^k}\}.
\]
这就得到 (iii)，并同时构造了局部积分叶。

实例。$\mathbb R^3$ 中的分布 $\mathcal D=\operatorname{span}\{\partial_x,-y\partial_z+\partial_y\}$：$[\partial_x,-y\partial_z+\partial_y]=0\in\mathcal D$，对合，积分叶为平面 $y=\text{const}$（每个叶是 $xz$ 平面）。反例 $\mathcal D=\operatorname{span}\{\partial_x,x\partial_z+\partial_y\}$：$[\partial_x,x\partial_z+\partial_y]=\partial_z\notin\mathcal D$，不对合，不存在通过每点的二维积分曲面——这是接触几何的典型例子。Frobenius 定理在控制论中给出可达性条件：向量场张成的分布对合当且仅当系统可沿积分叶任意移动。
\end{derivation}

\section{微分理想与 Pfaff 形式}

向量场版本的 involutivity 可以完全转写到余切空间。若 \(\alpha\in\Gamma(\mathcal D^\circ)\) 且 \(X,Y\in\Gamma(\mathcal D)\)，则
\[
\alpha(X)=\alpha(Y)=0.
\]
外微分公式给出
\begin{equation}\label{eq:math-dg-frobenius-dalpha-bracket}
\dd\alpha(X,Y)
=-\alpha([X,Y]).
\end{equation}
因此 \([X,Y]\in\mathcal D\) 当且仅当每个湮灭一形式的 \(\dd\alpha\) 在两个 \(\mathcal D\) 方向上都为零。在线性代数上，这等价于 \(\dd\alpha\) 至少含一个来自 \(\mathcal D^\circ\) 的因子，从而得到\eqnrefs{eq:math-dg-frobenius-pfaff}。

若余维为一，\(\mathcal D=\ker\alpha\) 且 \(\alpha\) 处处非零，Frobenius 条件进一步压缩成
\begin{equation}\label{eq:math-dg-frobenius-codim-one}
\boxed{\alpha\wedge\dd\alpha=0.}
\end{equation}

\begin{derivation}[余维一 Pfaff 判据，即\eqnrefs{eq:math-dg-frobenius-codim-one} 与 Frobenius 条件的等价]
余维一时分布 $\mathcal D=\ker\alpha$ 由单一一形式 $\alpha$ 确定。Frobenius 条件 $\dd\alpha=\eta\wedge\alpha$ 等价于 $\alpha\wedge\dd\alpha=0$，后者是可直接检验的代数条件。关键观察是：在 $n$ 维流形上，$\alpha\wedge\dd\alpha$ 是一个三形式，它在每点只张成一个一维空间（或为零），因此其消失条件只含 $\binom n3$ 个分量函数，远少于检验全部 Lie 括号。

方向一：Frobenius $\Rightarrow$ Pfaff。若 $\dd\alpha=\eta\wedge\alpha$，则
\begin{equation*}
 \alpha\wedge\dd\alpha
 =\alpha\wedge\eta\wedge\alpha=0,
\end{equation*}
因为 $\alpha\wedge\alpha=0$（一形式与自身的楔积为零）。这一方向不需要 $\eta$ 的任何具体性质，只要它存在即可。

方向二：Pfaff $\Rightarrow$ Frobenius。假设 $\alpha\wedge\dd\alpha=0$。局部取向量场 $N$ 使 $\alpha(N)=1$（归一化），并定义一形式
\begin{equation*}
 \eta:=-\iota_{N}\dd\alpha,
\end{equation*}
即 $\eta$ 是 $\dd\alpha$ 沿 $N$ 方向的缩并（带负号）。这里 $N$ 的存在性用到 $\alpha$ 处处非零，从而可在每点的余切空间中选一个对偶向量；局部光滑性由 $\alpha$ 光滑保证。

验证 $\dd\alpha=\eta\wedge\alpha$。需在基向量对上逐一检验。取 $X,Y\in\ker\alpha$，并把 $\{X,Y,N\}$ 当作局部标架（在余维一情形下 $\ker\alpha$ 处处补上 $N$ 即给出 $TM$）：
\begin{itemize}
\item 在 $(X,Y)$ 上：由 $\alpha\wedge\dd\alpha=0$ 在 $(N,X,Y)$ 上取值，展开
\begin{equation*}
(\alpha\wedge\dd\alpha)(N,X,Y)
=\alpha(N)\dd\alpha(X,Y)-\alpha(X)\dd\alpha(N,Y)+\alpha(Y)\dd\alpha(N,X),
\end{equation*}
后两项因 $\alpha(X)=\alpha(Y)=0$ 消失，只剩 $\alpha(N)\dd\alpha(X,Y)=\dd\alpha(X,Y)=0$。而 $(\eta\wedge\alpha)(X,Y)=\eta(X)\alpha(Y)-\eta(Y)\alpha(X)=0$。两式相等。
\item 在 $(X,N)$ 上：$(\eta\wedge\alpha)(X,N)=\eta(X)\alpha(N)-\eta(N)\alpha(X)=\eta(X)=-\dd\alpha(N,X)=\dd\alpha(X,N)$。两式相等。这里最后一步用了 $\eta$ 的定义与二形式的反对称性。
\item 在 $(N,N)$ 上：两二形式在重复向量上取值均为零。
\end{itemize}
在所有基向量对上取值相同，故
\begin{equation*}
 \dd\alpha=\eta\wedge\alpha.
\end{equation*}

积分流形的显式构造。在余维一可积情形下，Frobenius 定理给出局部坐标 $(x^1,\ldots,x^{n-1},u)$ 使
\begin{equation*}
\alpha=\mu(x,u)\,\dd u,
\end{equation*}
其中 $\mu$ 是处处非零的光滑函数。叶状结构由 $u=\text{const}$ 给出，每片叶是 $n-1$ 维子流形。等价地，存在局部积分因子 $\mu^{-1}$ 使 $\mu^{-1}\alpha$ 成为恰当形式 $\dd u$；这正是经典 Pfaff 方程 $\alpha=0$ 求解的第一积分存在性。

接触结构的反例。在 $\RR^3$ 上取 $\alpha=\dd z-x\,\dd y$，则 $\dd\alpha=-\dd x\wedge\dd y$，于是
\begin{equation*}
\alpha\wedge\dd\alpha
=-(\dd z-x\,\dd y)\wedge\dd x\wedge\dd y
=-\dd z\wedge\dd x\wedge\dd y\neq0.
\end{equation*}
Pfaff 判据立即给出 $\ker\alpha$ 不可积，不存在二维积分曲面。这正是接触几何的出发点：$\alpha\wedge\dd\alpha\neq0$ 处处成立时 $\alpha$ 称为接触形式，$\ker\alpha$ 称为接触分布，它"最大限度地不可积"。

叶状结构的正例。在 $\RR^3$ 上取 $\alpha=\dd z-x\,\dd x$，则 $\dd\alpha=-\dd x\wedge\dd x=0$，从而 $\alpha\wedge\dd\alpha=0$。由于 $\alpha=\dd(z-\tfrac12x^2)$ 是恰当形式，等值面 $z-\tfrac12x^2=c$ 的切空间正是 $\ker\alpha$，因而它们构成所求的抛物柱面叶状结构。

于是\eqnrefs{eq:math-dg-frobenius-codim-one}与 Frobenius 条件局部等价。Pfaff 判据 $\alpha\wedge\dd\alpha=0$ 的优势在于不需要显式构造 $\eta$，只需检查一个三形式的消失；在余维一情形下它把可积性问题完全化为一个代数条件，这是接触几何与叶状结构理论分歧的精确分水岭。
\end{derivation}

\begin{exbox}[接触型平面场不可积]
在 \(\RR^3\) 上取
\[
\alpha=\dd z-x\,\dd y,
\qquad
\mathcal D=\ker\alpha.
\]
有
\[
\dd\alpha=-\dd x\wedge\dd y,
\]
因此
\[
\alpha\wedge\dd\alpha
=-(\dd z-x\dd y)\wedge\dd x\wedge\dd y
=-\dd z\wedge\dd x\wedge\dd y\neq0.
\]
所以 \(\mathcal D\) 不可积：不存在处处以这些二维平面为切平面的曲面族。这一例子说明“每点都给出一个二维平面”远不足以保证它们能拼成曲面。
\end{exbox}


接下来考察接触结构、Reeb 场与 bracket-generating 分布。

余维一 Pfaff 条件的另一端不是“稍微不可积”，而是一个非常刚性的极大不可积结构。设 $M^{2m+1}$ 为奇数维流形，若一形式 $\alpha$ 满足
\begin{equation}
\boxed{
\alpha\wedge(\dd\alpha)^m\neq0
\quad\text{处处成立},
}
\label{eq:math-dg-contact-condition}
\end{equation}
则称 $\alpha$ 为 contact form，超平面分布
\[
\mathcal H:=\ker\alpha
\]
称为 contact distribution。条件\eqnrefs{eq:math-dg-contact-condition}立刻说明 $\mathcal H$ 不可能满足 Frobenius：若 $\dd\alpha=\eta\wedge\alpha$，则 $(\dd\alpha)^m$ 每一项都含 $\alpha$ 因子，从而 $\alpha\wedge(\dd\alpha)^m=0$，与接触条件矛盾。

更精确地，$\dd\alpha$ 在 $\mathcal H$ 上是非退化二形式。若 $X\in\mathcal H_p$ 且
\[
\dd\alpha_p(X,Y)=0
\qquad\forall Y\in\mathcal H_p,
\]
则由线性代数可推出 $X=0$。因此每个 contact hyperplane 自身携带一个 $2m$ 维 symplectic vector space 结构
\begin{equation}
\left(\mathcal H_p,\dd\alpha_p|_{\mathcal H_p}\right).
\label{eq:math-dg-contact-hyperplane-symplectic}
\end{equation}
这说明 contact geometry 可以看成“奇数维流形上随点变化的一族 symplectic 超平面”。

contact form 还唯一确定一个横向向量场，称为 Reeb vector field $R_\alpha$：
\begin{equation}
\boxed{
\alpha(R_\alpha)=1,
\qquad
\iota_{R_\alpha}\dd\alpha=0.
}
\label{eq:math-dg-reeb-definition}
\end{equation}
存在唯一性来自分解
\[
T_pM=\mathcal H_p\oplus\RR R_\alpha(p).
\]
第一条件固定横向分量，第二条件利用 $\dd\alpha|_{\mathcal H}$ 的非退化性消除任何可能的水平修正。

\begin{derivation}[contact distribution 为什么一步 Lie 括号就生成全部切空间]
对水平向量场 $X,Y\in\Gamma(\mathcal H)$，有 $\alpha(X)=\alpha(Y)=0$。由\eqnrefs{eq:math-dg-frobenius-dalpha-bracket}，
\begin{equation}
\alpha([X,Y])=-\dd\alpha(X,Y).
\label{eq:math-dg-contact-bracket-vertical}
\end{equation}
固定一点 $p$ 和任意非零 $X_p\in\mathcal H_p$。由于接触条件等价于 $\dd\alpha|_{\mathcal H_p}$ 非退化，作为 $\mathcal H_p$ 上的辛形式，对每个非零 $X_p$ 必存在 $Y_p\in\mathcal H_p$ 使
\begin{equation*}
\dd\alpha_p(X_p,Y_p)\neq0.
\end{equation*}
具体地，取 $\mathcal H_p$ 的一组辛基 $\{e_i,f_j\}$ 使 $\dd\alpha_p(e_i,f_j)=\delta_{ij}$，将 $X_p$ 展开为 $X_p=\sum_i(a_i e_i+b_i f_i)$，则选 $Y_p=\sum_i(-b_i e_i+a_i f_i)$ 即给出
\begin{equation*}
\dd\alpha_p(X_p,Y_p)=\sum_i(a_i^2+b_i^2)=\|X_p\|^2>0.
\end{equation*}
把 $X_p,Y_p$ 延拓为局部水平场 $X,Y$（即 $\alpha(X)=\alpha(Y)=0$ 在邻域内成立），便由\eqnrefs{eq:math-dg-contact-bracket-vertical}得到
\begin{equation*}
\alpha_p([X,Y]_p)=-\dd\alpha_p(X_p,Y_p)\neq0.
\end{equation*}

设 $R_\alpha$ 为 Reeb 向量场，满足 $\alpha(R_\alpha)=1$、$\iota_{R_\alpha}\dd\alpha=0$。由 $\alpha_p([X,Y]_p)\neq0$，可写
\begin{equation*}
[X,Y]_p
=\alpha_p([X,Y]_p)\,R_\alpha(p)+H_p,
\qquad H_p\in\mathcal H_p,
\end{equation*}
所以 $[X,Y]_p$ 具有非零 Reeb 方向分量。又 $\mathcal H_p=\ker\alpha_p$ 维数为 $2m$，而 $T_pM$ 维数为 $2m+1$，故
\begin{equation*}
\dim\bigl(\mathcal H_p+[\mathcal H,\mathcal H]_p\bigr)
\ge\dim\mathcal H_p+1
=2m+1
=\dim T_pM,
\end{equation*}
反向不等式自动成立，即
\begin{equation}
\boxed{
\mathcal H_p+[\mathcal H,\mathcal H]_p=T_pM.
}
\label{eq:math-dg-contact-step-two-bracket-generating}
\end{equation}
contact distribution 因而是 step-two bracket-generating distribution：虽然单个允许速度永远限制在 $\mathcal H$ 中，但通过"先沿 $X$、再沿 $Y$、再反向返回"的小回路，可以在二阶产生横向位移。
\end{derivation}

这个现象与 Frobenius 恰好相反。Frobenius 条件要求 Lie brackets 不产生新方向，于是运动被困在低维叶片中；bracket-generating 条件则要求反复 Lie brackets 最终张成整个切空间。Chow--Rashevskii 定理说明，在连通流形上若一个光滑分布 bracket generating，则任意两点都可由分段光滑的水平曲线连接。若再在 $\mathcal H$ 上给定内积，便得到 sub-Riemannian geometry：距离只计算水平曲线长度，而“缺失方向”由非交换运动间接产生。

最基本的局部标准型由 contact Darboux theorem 给出。任一点附近都存在坐标
\[
(q^1,\ldots,q^m,p_1,\ldots,p_m,z)
\]
使 contact form 在乘以一个处处非零函数后可写成
\begin{equation}
\alpha
=\dd z-\sum_{a=1}^m p_a\,\dd q^a.
\label{eq:math-dg-contact-darboux-form}
\end{equation}
因此 contact geometry 与 symplectic geometry 一样没有局部标量曲率型不变量：所有 contact structures 在足够小邻域内都有相同标准形，真正的信息主要来自全局拓扑、Reeb dynamics 与附加结构。

contact form 还可以通过增加一个实坐标变成真正的 symplectic form。令
\[
\widetilde M:=\RR_s\times M,
\qquad
\lambda:=\ee^s\alpha,
\]
并定义
\begin{equation}
\boxed{
\omega:=\dd\lambda
=\ee^s(\dd s\wedge\alpha+\dd\alpha).
}
\label{eq:math-dg-contact-symplectization}
\end{equation}
由 $\dd^2=0$ 有 $\dd\omega=0$。更重要的是
\begin{align*}
\omega^{m+1}
&=(m+1)\ee^{(m+1)s}
\dd s\wedge\alpha\wedge(\dd\alpha)^m,
\end{align*}
由 contact 条件处处非零。因此 $\omega$ 非退化，$(\RR\times M,\omega)$ 是 symplectic manifold，称为 $(M,\alpha)$ 的 symplectization。contact geometry 因而不是与 symplectic geometry 平行的孤立理论，而是奇数维边界数据在增加一个尺度方向后的偶数维 symplectic 几何。

向量场 $\partial_s$ 是这个 exact symplectic manifold 的 Liouville vector field，因为
\begin{equation}
\iota_{\partial_s}\omega
=\ee^s\alpha=\lambda,
\qquad
\mathcal L_{\partial_s}\omega=\omega.
\label{eq:math-dg-symplectization-liouville}
\end{equation}
所以沿 $s$ 方向平移会把 symplectic form 作指数缩放。另一方面，Reeb 场在 symplectization 中恰好由径向 Hamiltonian 识别：若
\[
H(s,x)=\ee^s,
\]
并采用 Hamiltonian 约定
\[
\iota_{X_H}\omega=-\dd H,
\]
则
\begin{equation}
X_H=R_\alpha.
\label{eq:math-dg-reeb-as-hamiltonian}
\end{equation}
因为
\[
\iota_{R_\alpha}\omega
=\ee^s\iota_{R_\alpha}(\dd s\wedge\alpha+\dd\alpha)
=-\ee^s\dd s
=-\dd(\ee^s).
\]
这给 Reeb dynamics 一个非常具体的解释：它是在 symplectization 上由纯尺度 Hamiltonian 产生的 Hamilton flow。

在 contact manifold 本身，也可为任意光滑函数 $H\in C^\infty(M)$ 定义 contact Hamiltonian vector field $X_H$。固定约定
\begin{equation}
\boxed{
\alpha(X_H)=H,
\qquad
\iota_{X_H}\dd\alpha
=(R_\alpha H)\alpha-\dd H.
}
\label{eq:math-dg-contact-hamiltonian}
\end{equation}
第二式在 $\mathcal H$ 上唯一确定 $X_H$ 的水平分量，第一式再确定 Reeb 分量。由 Cartan 公式立刻得到
\begin{equation}
\boxed{
\mathcal L_{X_H}\alpha
=(R_\alpha H)\alpha.
}
\label{eq:math-dg-contact-hamiltonian-conformal}
\end{equation}
因此 contact Hamiltonian flow 一般不保持具体的一形式 $\alpha$，而保持其 kernel 所定义的 contact distribution；当 $R_\alpha H=0$ 时才严格保持 $\alpha$。

\begin{derivation}[contact Hamiltonian 如何提升为 symplectic Hamiltonian]
在 symplectization 上令
\[
\widetilde H(s,x)=\ee^sH(x).
\]
把 Hamiltonian vector field 分解为
\[
X_{\widetilde H}=X_H+a\,\partial_s.
\]
利用\eqnrefs{eq:math-dg-contact-symplectization}，逐项计算
\[
\iota_{X_{\widetilde H}}\omega
=\ee^s\left[
a\alpha-H\dd s+\iota_{X_H}\dd\alpha
\right].
\]
而
\[
-\dd\widetilde H
=-\ee^s(H\dd s+\dd H).
\]
比较两式并代入\eqnrefs{eq:math-dg-contact-hamiltonian}，得到
\[
a=-R_\alpha H.
\]
所以
\begin{equation}
\boxed{
X_{\widetilde H}
=X_H-(R_\alpha H)\partial_s.
}
\label{eq:math-dg-contact-hamiltonian-lift}
\end{equation}

齐次性的显式验证。对常数 $\lambda>0$，记 $r_\lambda(s,x)=(s+\ln\lambda,x)$ 为沿 $s$ 方向的平移。由 $\widetilde H$ 的定义，
\begin{equation*}
\widetilde H\circ r_\lambda
=\ee^{s+\ln\lambda}H(x)
=\lambda\,\widetilde H(s,x),
\end{equation*}
即 $\widetilde H$ 是沿 $s$ 方向的 $1$ 次齐次函数。对应地，symplectic 形式 $\omega=\dd(\ee^s\alpha)$ 在 $r_\lambda$ 下满足 $r_\lambda^*\omega=\lambda\,\omega$，因此 Hamiltonian 流保持这一齐次结构：若 $\phi_t$ 是 $\widetilde H$ 的流，则 $r_\lambda\circ\phi_t\circ r_{\lambda^{-1}}$ 也是 $\widetilde H$ 的流，由唯一性得 $r_\lambda\circ\phi_t=\phi_t\circ r_\lambda$。

Reeb 流作为特例。取 $H\equiv1$，则 $R_\alpha H=0$ 且 $X_H=R_\alpha$（由 contact Hamiltonian 的定义方程 $\iota_{X_H}\dd\alpha=\dd H-(R_\alpha H)\alpha=0$ 与 $\alpha(X_H)=H=1$）。此时
\begin{equation*}
X_{\widetilde H}=R_\alpha,
\end{equation*}
即 Reeb 流恰好提升为 symplectization 上与 $s$ 坐标无关的 Hamiltonian 流。这给出 Reeb 流所有接触不变量的 symplectic 解释。

三维接触情形的显式公式。在标准接触 $\RR^3$ 上取 $\alpha=\dd z-\tfrac12(x\,\dd y-y\,\dd x)$，Reeb 场 $R_\alpha=\partial_z$。对 $H=H(x,y,z)$，contact Hamiltonian 场为
\begin{equation*}
X_H
=\left(\frac{\partial H}{\partial y}+\frac{x}{2}\frac{\partial H}{\partial z}\right)\partial_x
-\left(\frac{\partial H}{\partial x}-\frac{y}{2}\frac{\partial H}{\partial z}\right)\partial_y
+\left(x\frac{\partial H}{\partial y}-y\frac{\partial H}{\partial x}-H\right)\partial_z,
\end{equation*}
提升到 $\RR^4_{(s,x,y,z)}$ 后沿 $s$ 方向的分量为 $-\partial_z H$。当 $H$ 与 $z$ 无关时提升为齐次 Hamiltonian 系统，这正是经典 Hamilton 力学在接触框架下的对应。

物理建模：接触 Hamilton 系统。经典耗散系统的 Hamilton 化需要接触几何。对带摩擦的一维振子 $\ddot q+\gamma\dot q+\omega^2q=0$，取接触形式 $\alpha=\dd z-p\,\dd q$ 与 Hamiltonian $H=\tfrac12p^2+\tfrac12\omega^2q^2+\gamma z$，则 $R_\alpha H=\gamma$，contact Hamiltonian 流给出原耗散方程，而能量沿流的变化为
\begin{equation*}
\frac{\dd H}{\dd t}=R_\alpha H\cdot H=\gamma H,
\end{equation*}
正确反映能量耗散（或注入）。提升到 symplectization 后，$\widetilde H=\ee^sH$ 的 symplectic 流在 $s$ 方向以速率 $-\gamma$ 演化，把耗散编码为 symplectic 框架下的额外自由度。

contact Hamiltonian dynamics 因而可以看成一类对 $s$ 方向齐次的 symplectic Hamiltonian dynamics；Reeb flow 对应 $H\equiv1$ 的特例，而一般情形下 Reeb 分量 $-R_\alpha H$ 正好度量系统偏离保守的程度。
\end{derivation}

\begin{exbox}[三维 Heisenberg contact structure]
在 $\RR^3_{(x,y,z)}$ 上取
\[
\alpha=\dd z-\frac12(x\,\dd y-y\,\dd x).
\]
则
\[
\dd\alpha=-\dd x\wedge\dd y,
\qquad
\alpha\wedge\dd\alpha=-\dd z\wedge\dd x\wedge\dd y\neq0.
\]
一组水平向量场可取
\[
X=\partial_x-\frac y2\partial_z,
\qquad
Y=\partial_y+\frac x2\partial_z,
\]
而 Reeb 场为 $R_\alpha=\partial_z$。直接计算
\[
[X,Y]=\partial_z=R_\alpha.
\]
所以被禁止的竖直方向恰由两个水平运动的交换子生成。这是 Heisenberg group、非完整约束与 sub-Riemannian 控制理论共同的局部模型。
\end{exbox}

\section{从局部可积到全局叶空间：holonomy 与主丛水平分布}

Frobenius 定理是局部定理。若 \(\mathcal D\) involutive，通过每一点的局部积分流形可以沿重叠区域唯一延拓，并得到一个极大的连通积分叶；这些极大叶组成 foliation \(\mathcal F\)。但“每点都有局部叶”并不保证商空间
\[
 M/\mathcal F
\]
本身是流形，甚至不保证 Hausdorff。全局问题取决于叶的闭性、holonomy 以及叶在 \(M\) 中的回返方式，而这些信息都不包含在点态条件 \([X,Y]\in\Gamma(\mathcal D)\) 中。

最简单的反例是二维环面上的 Kronecker foliation。取
\[
 M=\mathbb T^2=\mathbb R^2/(2\pi\mathbb Z)^2,
 \qquad
 X=\partial_x+\alpha\partial_y,
 \qquad
 \alpha\notin\mathbb Q,
\]
并令 \(\mathcal D=\operatorname{span}\{X\}\)。秩一分布自动 involutive，因此每条积分曲线
\[
 t\longmapsto
 (x_0+t,\,y_0+\alpha t)\pmod{2\pi}
\]
都是一片叶。由于 \(\alpha\) 无理，每片叶在 \(\mathbb T^2\) 中稠密。于是叶空间中的单点原像并不是闭集，商空间甚至不是 \(T_1\)，因而不可能是 Hausdorff 流形。这个例子把“局部 Frobenius 坐标存在”和“存在良好的全局叶坐标”严格区分开。

“叶空间不一定良好”并不意味着任何全局纤维化都不可能。一个重要充分条件由 Ehresmann fibration theorem 给出。

\begin{theorem}[Ehresmann fibration theorem]\label{thm:math-dg-ehresmann-fibration}
设 $F:M\to B$ 是光滑、满射、proper submersion。则 $F$ 是局部平凡光滑纤维丛：对每个 $b_0\in B$，存在邻域 $U$ 和微分同胚
\begin{equation}
F^{-1}(U)\cong U\times F^{-1}(b_0)
\label{eq:math-dg-ehresmann-local-trivialization}
\end{equation}
使 $F$ 对应第一坐标投影。
\end{theorem}

submersion 条件保证每个 fiber 都是同维光滑子流形；properness 则提供把局部水平提升延拓到整个小邻域所需的全局控制。取任意 Riemannian metric，并令
\[
\mathcal V=\ker\dd F,
\qquad
\mathcal H=\mathcal V^\perp.
\]
对 $B$ 上的向量场 $Y$，$\dd F|_{\mathcal H}$ 逐点是同构，所以存在唯一水平提升 $Y^H$。若没有 properness，这些水平提升可能在有限时间逃向无穷远，使不同 fiber 无法由流一致识别；properness 恰好排除这一失控。

\begin{derivation}[Ehresmann 定理的局部平凡化主线]
在 $b_0$ 附近选坐标球 $U$ 使其闭包紧，并对每个 $b\in U$ 取从 $b_0$ 指向 $b$ 的径向路径 $\gamma_b:[0,1]\to U$，满足 $\gamma_b(0)=b_0$、$\gamma_b(1)=b$。把 $\dot\gamma_b(t)$ 水平提升到 $M$，从任意 $x\in F^{-1}(b_0)$ 出发解 ODE
\[
\dot{\widetilde\gamma}_{b,x}(t)
=\dot\gamma_b(t)^H,
\qquad
\widetilde\gamma_{b,x}(0)=x.
\]
沿提升曲线计算 $F$ 的变化率，由链式法则与水平提升的定义，
\begin{align*}
\frac{\dd}{\dd t}F(\widetilde\gamma_{b,x}(t))
&=\dd F_{\widetilde\gamma_{b,x}(t)}\bigl(\dot{\widetilde\gamma}_{b,x}(t)\bigr)\\
&=\dd F_{\widetilde\gamma_{b,x}(t)}\bigl(\dot\gamma_b(t)^H\bigr)
=\dot\gamma_b(t),
\end{align*}
最后一步用到 $\dd F|_{\mathcal H}$ 是恒等提升。积分两端并代入初值 $F(\widetilde\gamma_{b,x}(0))=F(x)=b_0=\gamma_b(0)$，得到
\[
F(\widetilde\gamma_{b,x}(t))=\gamma_b(t),
\]
因此提升曲线始终停留在对应 fiber 之上。因为 $F$ proper，紧小球 $\overline U$ 的原像 $F^{-1}(\overline U)$ 紧，水平 ODE 的解不会在有限时间逃向无穷远；缩小 $U$ 后，所有这些 ODE 在 $t\in[0,1]$ 上统一存在。定义
\[
\Phi(b,x):=\widetilde\gamma_{b,x}(1).
\]
由 ODE 对初值 $x$ 与参数 $b$ 的光滑依赖，$\Phi$ 光滑；且由上式 $F(\Phi(b,x))=\gamma_b(1)=b$。构造逆映射时，对 $(y,b)\in F^{-1}(U)$ 取从 $b$ 沿径向返回 $b_0$ 的路径 $\gamma_b(1-t)$，从 $y$ 出发水平提升至 $t=1$，得到 $F^{-1}(b_0)$ 中的一点，与 $b$ 组合即给出 $\Phi^{-1}$，它同样由 ODE 光滑依赖保证光滑。因此
\[
\Phi:U\times F^{-1}(b_0)\to F^{-1}(U)
\]
是微分同胚，并满足 $F\circ\Phi(b,x)=b$，这正得到\eqnrefs{eq:math-dg-ehresmann-local-trivialization}。

properness 的必要性。考虑 $F:\RR^2\setminus\{0\}\to\RR$，$F(x,y)=x$。这是 submersion 但不 proper：fiber $F^{-1}(0)=\{(0,y):y\neq0\}$ 不紧。取 $b_0=0$、$b_n\to0^+$，从 $(0,1)$ 出发沿 $b_n$ 方向水平提升，当 $b_n\to0$ 时提升曲线可能逃向无穷远，不存在统一局部平凡化。properness 通过紧性排除这种逃逸。

纤维丛结构的进一步推论。当 $F:M\to B$ 是 proper submersion 且 $B$ 连通时，所有 fiber 微分同胚。结合 long exact sequence of homotopy groups，
\begin{equation*}
\cdots\to\pi_1(F)\to\pi_1(M)\to\pi_1(B)\to\pi_0(F)\to\cdots
\end{equation*}
给出纤维丛的同伦分类。Ehresmann 定理因此是 fibrations 在光滑范畴的版本：proper submersion 自动是 locally trivial fibration。

举例：Hopf 纤维化。取 $S^3\subset\RR^4$ 为单位三维球，$S^2=\CP^1$ 为 Riemann 球，Hopf 映射 $F:S^3\to S^2$ 由
\begin{equation*}
F(z_1,z_2)=\bigl(2\operatorname{Re}(z_1\bar z_2),\,2\operatorname{Im}(z_1\bar z_2),\,|z_1|^2-|z_2|^2\bigr)
\end{equation*}
给出。$S^3$ 紧、$F$ submersion，故自动 proper，Ehresmann 给出 $S^3\to S^2$ 是 $S^1$-纤维丛。水平分布对应 Hopf 联络，其曲率是 $S^2$ 上面积形式的拉回，正是 Berry 几何相的数学来源。

与叶状结构的对比。Ehresmann 的纤维丛结构是"被动"的：从一个 proper submersion 推出纤维自动形成局部平凡族。叶状结构理论则是"主动"的：从一个可积分布出发构造叶，但叶空间可能具有非流形拓扑（例如 Reeb 叶状结构中叶空间含区间与 Cantor 集混合）。两者代表可积性在全局层面的两种不同表现。

光滑依赖的精确陈述。ODE 对初值与参数的光滑依赖定理要求向量场关于参数光滑、关于初值 Lipschitz。这里水平提升 $\dot\gamma_b(t)^H$ 关于 $b$ 光滑（因 $\dd F|_\mathcal H$ 光滑可逆）、关于 $(t,x)$ 光滑（因水平分布光滑），因此 $\Phi(b,x)$ 是 $C^\infty$ 的。逆映射 $\Phi^{-1}$ 由反向 ODE 给出，同样光滑。这是 Ehresmann 定理把 ODE 理论与微分拓扑衔接的关键点。

一般结论：Ehresmann 纤维化定理。若 $F:M\to B$ 是 proper submersion，则 $F$ 是 locally trivial fibration。结合 $B$ 连通条件，所有 fiber 微分同胚；进一步若 $B$ 单连通，纤维丛的整体结构由任一 fiber 的同伦型决定。这把纤维丛分类问题化为同伦论问题，是 characteristic class 理论的出发点。Ehresmann 定理因此是微分拓扑中"properness 换全局平凡性"的典型范式，与 Sard 定理、横截性定理同属 Morse 理论的工具族。

物理建模：Berry 联络与 Ehresmann。在量子绝热系统中，参数空间 $B$ 上的 Hilbert 丛 $\mathcal H\to B$ 的投影是 submersion，Ehresmann 给出局部平凡化。Berry 联络是水平分布的选择，曲率即 Berry 曲率。当参数空间含简并点时 properness 可能失效，对应量子相变，这是 Ehresmann 条件在物理上的具体意义。
\end{derivation}

所以 Frobenius、Ehresmann 与 holonomy 处在不同层次：Frobenius 判断一个给定分布是否局部来自叶；Ehresmann 从一个 proper submersion 推出 fibers 全局形成局部平凡族；holonomy 则记录沿闭路运输后 fiber 数据如何回返。三者不能互相替代。

主丛联络给出一个更结构化的全局版本。设
\[
 \pi:P\to B
\]
是主 \(G\)-丛，\(\mathfrak g\) 为 Lie 代数。垂直分布是
\[
 \mathcal V=\ker\dd\pi,
\]
而主联络一形式 \(\omega\in\Omega^1(P,\mathfrak g)\) 定义水平分布
\begin{equation}\label{eq:math-dg-principal-horizontal-distribution}
 \mathcal H:=\ker\omega,
 \qquad
 TP=\mathcal H\oplus\mathcal V.
\end{equation}
曲率二形式为
\begin{equation}\label{eq:math-dg-principal-curvature}
 \Omega
 =\dd\omega+\frac12[\omega\wedge\omega].
\end{equation}
若 \(X^H,Y^H\) 是水平向量场，则 \(\omega(X^H)=\omega(Y^H)=0\)，故
\begin{equation}\label{eq:math-dg-curvature-nonintegrability}
 \boxed{
 \Omega(X^H,Y^H)
 =-\omega([X^H,Y^H]).}
\end{equation}
因此
\begin{equation}\label{eq:math-dg-horizontal-frobenius}
 \boxed{
 \mathcal H\text{ involutive}
 \quad\Longleftrightarrow\quad
 \Omega=0.}
\end{equation}
这里的等价使用了主联络曲率本来就是 horizontal 的事实。曲率不再只是“平行移动绕小回路后的误差”，它就是水平分布违反 Frobenius 条件的精确张量。

\begin{derivation}[曲率为何等于水平分布的 Frobenius 障碍]
对任意 \(\mathfrak g\)-值一形式 \(\omega\) 与任意向量场 \(X,Y\)，外微分的 invariant 公式为
\[
 \dd\omega(X,Y)
 =X\{\omega(Y)\}-Y\{\omega(X)\}-\omega([X,Y]).
\]
令 \(X=X^H\)、\(Y=Y^H\) 为水平提升。由水平分布定义 \(\mathcal H=\ker\omega\)，有 \(\omega(X^H)=\omega(Y^H)=0\)，前两项均为零，剩下
\[
 \dd\omega(X^H,Y^H)=-\omega([X^H,Y^H]).
\]
再算 wedge-bracket 项。由定义
\[
 [\omega\wedge\omega](X^H,Y^H)
 =[\omega(X^H),\omega(Y^H)]-[\omega(Y^H),\omega(X^H)],
\]
每个因子均为零，故
\[
 [\omega\wedge\omega](X^H,Y^H)=0.
\]
把两式代入曲率定义 \(\Omega=\dd\omega+\tfrac12[\omega\wedge\omega]\)，得到
\begin{align*}
\Omega(X^H,Y^H)
&=\dd\omega(X^H,Y^H)+\tfrac12[\omega\wedge\omega](X^H,Y^H)\\
&=-\omega([X^H,Y^H])+0,
\end{align*}
即\eqnrefs{eq:math-dg-curvature-nonintegrability}。若 \(\Omega=0\)，则对一切水平对有 \(\omega([X^H,Y^H])=0\)，即 \([X^H,Y^H]\in\ker\omega=\mathcal H\)，括号仍水平，Frobenius 定理给出局部水平叶。反之若水平分布 involutive，则 \([X^H,Y^H]\in\mathcal H\)，故 \(\omega([X^H,Y^H])=0\)，由上式曲率在全部水平对上消失；主联络曲率又满足 horizontality，即它被水平向量场完全决定，因此 \(\Omega=0\) 处处成立，即\eqnrefs{eq:math-dg-horizontal-frobenius}。

Bianchi 恒等式的 Frobenius 解释。第一 Bianchi 恒等式 $D\Omega=0$ 可视为水平分布的"高阶可积性"约束：若 $\Omega=0$ 给出局部叶，则 $D\Omega=0$ 自动满足；若 $\Omega\neq0$，Bianchi 限制曲率沿水平方向变化的可能方式。第二 Bianchi 恒等式 $D^2=0$ 则来自外微分的闭性 $\dd^2=0$，与 Frobenius 无直接关联，而是曲率作为二形式的代数性质。

Abelian 与非 Abelian 情形的差异。当结构群 $G$ 为 Abelian（如 $U(1)$ 电磁联络），$[\omega\wedge\omega]=0$，曲率简化为 $\Omega=\dd\omega$。此时 Frobenius 障碍就是 $\dd\omega=0$，即 $\omega$ 闭；局部叶存在当且仅当 $\omega$ 局部恰当，这正是电磁势的规范不变性 $\omega\mapsto\omega+\dd f$ 不改变曲率。非 Abelian 情形下 $[\omega\wedge\omega]$ 项给出非交换性贡献，使曲率不再是线性算子，Yang--Mills 理论的非线性性正源于此。

holonomy 的具体计算。对 $U(1)$ 联络在圆 $S^1$ 上的平行移动，沿闭路 $\gamma$ 的 holonomy 为
\begin{equation*}
\operatorname{Hol}(\gamma)=\exp\left(i\oint_\gamma\omega\right)
=\exp\left(i\int_D\Omega\right),
\end{equation*}
最后一步用到 Stokes 定理。曲率 $\Omega$ 因此是 holonomy 的"密度"：把闭路包围的曲率积分即得 holonomy 角度。非 Abelian 情形下此式变为非交换 Wilson 圈
\begin{equation*}
W(\gamma)=\mathcal P\exp\left(\oint_\gamma\omega\right),
\end{equation*}
其中 $\mathcal P$ 表示 path-ordering，反映非交换性使局部贡献不能简单交换。

但 \(\Omega=0\) 仍只是局部可积性。固定 \(p_0\in P\) 后，沿基流形路径作水平提升定义平行移动。若基空间存在非平凡闭路，平坦联络仍可能具有非平凡 holonomy representation
\[
 \operatorname{Hol}:\pi_1(B,b_0)\to G.
\]
只有当这些闭路 holonomy 全部平凡时，从 \(p_0\) 出发的水平叶才会在每个基点上选出唯一一点，从而成为全局水平截面。若 \(B\) 单连通，则 \(\pi_1(B)=0\)，平坦性使平行移动与路径无关，于是该局部 Frobenius 结构可以沿整个连通基底拼成全局水平截面。由此可见：曲率控制局部可积性，holonomy 控制全局拼接；前者是微分条件，后者是拓扑条件，两者共同决定联络的可积性结构。

分类空间与特征类。曲率与 holonomy 的分离对应特征类的两层结构：第一 Chern 类 $c_1\in H^2(B,\ZZ)$ 通过 de Rham 表示由曲率积分给出（局部信息），但其整体挠部分由 holonomy 决定（全局信息）。Chern--Weil 理论把曲率不变多项式映射到上同调类，给出 $c_k$、$p_k$、Euler 类等的微分代表；而 holonomy 的离散部分则对应 torsion 类，需要其他工具（如 $\hat A$-亏格、指数定理）才能完全捕捉。
\end{derivation}


接下来考察Holonomy group、曲率生成元与 Ambrose--Singer。

固定主丛点 $p\in P$，所有以 $b=\pi(p)$ 为基点的闭路通过水平提升给出结构群元素；这些元素组成 holonomy group
\[
\operatorname{Hol}_p(\omega)\subset G.
\]
其中由可缩闭路生成的连通部分称为 restricted holonomy，记作
\[
\operatorname{Hol}_p^0(\omega).
\]
平坦性 $\Omega=0$ 只保证 restricted holonomy 平凡；若基空间非单连通，完整 holonomy 仍可能来自基本群的离散表示。因此局部曲率与全局 monodromy 必须分开。

曲率为什么会生成 holonomy？考虑基流形上一块由向量 $X,Y$ 张成的无穷小平行四边形，沿其边界作平行移动。若面积尺度为 $\varepsilon^2$，则回到起点后的群元素具有展开
\[
\operatorname{Hol}_{\partial\square_\varepsilon}
=I-\varepsilon^2\Omega(X^H,Y^H)+O(\varepsilon^3),
\]
其中把 Lie algebra 元素通过相应表示解释为无穷小群作用。于是曲率正是 holonomy 的二阶局部生成元。

Ambrose--Singer 定理把这一局部图像升级为完整 Lie 代数结论。记 $\mathfrak{hol}_p$ 为 restricted holonomy group 的 Lie algebra，则
\begin{equation}
\boxed{
\mathfrak{hol}_p
=\operatorname{span}
\left\{
\operatorname{Ad}_{g_\gamma^{-1}}
\Omega_q(X^H,Y^H)
\right\}
}
\label{eq:math-dg-ambrose-singer}
\end{equation}
其中 $q$ 遍历可由从 $p$ 出发的水平路径到达的点，$g_\gamma$ 表示沿路径把曲率值平行搬回 $p$ 的群元素。换言之，不是在一个点读取 $\Omega_p$ 就能得到全部 holonomy；必须把各点曲率沿路径搬回同一 fiber 后再取 Lie algebra span。

对 Levi--Civita 联络，这一结论把 Riemann 曲率与切空间中的 holonomy group 联系起来。若 holonomy 严格小于 $SO(n)$，就意味着存在被平行移动保持的额外张量。例如 Kähler 几何中的复结构 $J$ 满足 $\nabla J=0$，因此 holonomy 降到 $U(m)$；Calabi--Yau 情形进一步降到 $SU(m)$。这里不展开特殊 holonomy 分类，但逻辑已经清楚：
\[
\text{曲率}
\Longrightarrow
\text{holonomy Lie algebra}
\Longrightarrow
\text{平行张量与额外几何结构}.
\]

\begin{derivation}[无穷小回路如何读出曲率]
取局部联络形式 $A=A_i\,\dd x^i$，沿 $X$、$Y$ 方向各走长度 $\varepsilon$ 的小矩形。平行移动在每条边上由 path-ordered exponential 给出。保留到二阶，四条边的乘积为
\[
\begin{aligned}
U_{\partial\square}
={}&
(I-\varepsilon A_X)
\bigl[I-\varepsilon(A_Y+\varepsilon\partial_XA_Y)\bigr]\\
&\times
\bigl[I+\varepsilon(A_X+\varepsilon\partial_YA_X)\bigr]
(I+\varepsilon A_Y)
+O(\varepsilon^3).
\end{aligned}
\]
一次项全部消失，二次项整理为
\[
U_{\partial\square}
=I-\varepsilon^2
\left(
\partial_XA_Y-
\partial_YA_X+[A_X,A_Y]
\right)
+O(\varepsilon^3).
\]
括号正是曲率分量
\[
F(X,Y)
=(\dd A+A\wedge A)(X,Y).
\]
因此曲率不是平行移动误差的比喻，而是无穷小闭路 holonomy 的首个非平凡系数。

逐项展开的细节。把四条边的乘积按 $\varepsilon$ 阶数展开。零阶项为 $I$。一阶项为
\begin{equation*}
-\varepsilon A_X-\varepsilon A_Y+\varepsilon A_X+\varepsilon A_Y=0,
\end{equation*}
正反向贡献相消，反映沿闭路回到起点的最低阶必然平凡。二阶项分三类：
\begin{itemize}
\item 同向导数项：$\varepsilon^2\partial_XA_Y$ 与 $-\varepsilon^2\partial_YA_X$，来自沿一边走时联络值在另一边方向的变化；
\item 交换项：$\varepsilon^2[A_X,A_Y]$，来自非交换 Lie 代数的乘积顺序差异；
\item 同边往返项：如 $(-\varepsilon A_X)(\varepsilon A_X)=-\varepsilon^2A_X^2$，但这些项在同一条边往返时已被 path-ordering 的局部解所吸收，不出现在最终公式中。
\end{itemize}
三类贡献合并即得曲率分量。

Abelian 情形的简化。当 $G$ 为 Abelian（如 $U(1)$），$[A_X,A_Y]=0$，曲率简化为 $F_{XY}=\partial_XA_Y-\partial_YA_X$，这正是电磁场强 $F_{\mu\nu}=\partial_\mu A_\nu-\partial_\nu A_\mu$。沿小回路的 holonomy 为
\begin{equation*}
U_{\partial\square}=\exp\left(-i\varepsilon^2F_{XY}\right)+O(\varepsilon^3),
\end{equation*}
对应磁通量等于回路包围的磁场面积元，这是 Faraday 定律的几何表述。

非 Abelian 情形与 Yang--Mills。对 $SU(N)$ 联络，交换项 $[A_X,A_Y]$ 给出非线性自相互作用，是 Yang--Mills 方程
\begin{equation*}
D^\mu F_{\mu\nu}=0
\end{equation*}
非线性的来源。物理上这对应胶子自相互作用：与电磁场光子不带电荷不同，胶子本身带色荷，因此曲率含 $A\wedge A$ 项。

Ambrose--Singer 定理的极限形式。把无穷小回路公式积分到有限回路，得到 holonomy 沿任意闭路 $\gamma$ 的表达式
\begin{equation*}
\operatorname{Hol}(\gamma)=\mathcal P\exp\left(-\oint_\gamma A\right),
\end{equation*}
对边界为 $\gamma$ 的曲面 $\Sigma$ 应用非交换 Stokes 定理，得
\begin{equation*}
\operatorname{Hol}(\gamma)=\mathcal P\mathcal S\exp\left(-\int_\Sigma F\right),
\end{equation*}
其中 $\mathcal S$ 表示 surface-ordering。把 $\Sigma$ 分成 $N^2$ 个面积为 $1/N^2$ 的小方格并取 $N\to\infty$，每个小方格贡献 $I-\frac1{N^2}F+O(N^{-3})$，乘积收敛到上式。这是 Ambrose--Singer 定理的显式版本：holonomy 由曲率沿曲面取值的非交换积分完全决定。

物理建模：Berry 相位与 Aharonov--Bohm。在量子力学中，参数空间上的绝热演化给出 $U(1)$ 联络（Berry 联络）$\mathcal A=i\langle\psi|\dd\psi\rangle$，曲率 $\mathcal F=\dd\mathcal A$ 是 Berry 曲率。沿参数空间闭路演化一周后波函数获得相位 $\exp(i\oint\mathcal A)$，这正是无穷小回路公式在 $U(1)$ 情形的应用。Aharonov--Bohm 效应中电子绕磁通量管的相位差则是非平凡 holonomy 的典型例子：联络局部平坦（磁场为零区域）但全局 holonomy 非平凡，体现曲率与 holonomy 的全局-局部分离。
\end{derivation}


Frobenius 定理因此应当分成两个层次理解：局部层次由 Lie 括号或微分理想决定，保证适配坐标和积分叶存在；全局层次还要检查叶空间拓扑与 holonomy。下一章把几何结构允许显含时间，研究张量、度量与由度量决定的联络在演化中的精确变化律。
